\documentclass[conference]{IEEEtran}
\IEEEoverridecommandlockouts
% The preceding line is only needed to identify funding in the first footnote. If that is unneeded, please comment it out.
\usepackage{cite}
\usepackage{longtable}
\usepackage{booktabs}
\usepackage{multirow}
\usepackage{amsmath,amssymb,amsfonts}
\usepackage{algorithmic}
\usepackage{graphicx}
\usepackage{textcomp}
\usepackage{xcolor}
\usepackage{enumitem}
\usepackage[listings, many]{tcolorbox}
\newcommand{\fm}[1]{\textcolor{magenta}{FM: #1}}
\newcommand{\sms}[1]{\textcolor{cyan}{SMS: #1}}
\newcommand{\method}{{LLM-TSFS}}
\def\BibTeX{{\rm B\kern-.05em{\sc i\kern-.025em b}\kern-.08em
    T\kern-.1667em\lower.7ex\hbox{E}\kern-.125emX}}
\begin{document}

\title{LLM-guided Feature Selection in Time Series Models for Small Data: the case of Migration Flows\\
{}
\thanks{This material is based upon work supported in part by the National Science Foundation for REU Site Grant 2349370.}
}

\author{
\IEEEauthorblockN{
  Phong Cao\IEEEauthorrefmark{1},
  Jacob Molnia\IEEEauthorrefmark{1},
  Elizabeth Stanish\IEEEauthorrefmark{2},
  Shannon Song,\IEEEauthorrefmark{1},
  Fabricio Murai\IEEEauthorrefmark{1}
}
\IEEEauthorblockA{
\IEEEauthorrefmark{1}Data Science Program, Worcester Polytechnic Institute, Worcester, MA, USA \\
\IEEEauthorrefmark{2}Computer Science Department,
Yale University, New Haven, CT, USA}
\IEEEauthorblockA{ptcao@wpi.edu, jrmolnia@wpi.edu, elizabeth.stanish@yale.edu, smsong@wpi.edu, fmurai@wpi.edu}
}

% \author{\IEEEauthorblockN{1\textsuperscript{st} Given Name Surname}
% \IEEEauthorblockA{\textit{dept. name of organization (of Aff.)} \\
% \textit{name of organization (of Aff.)}\\
% City, Country \\
% email address or ORCID}
% \and
% \IEEEauthorblockN{2\textsuperscript{nd} Given Name Surname}
% \IEEEauthorblockA{\textit{dept. name of organization (of Aff.)} \\
% \textit{name of organization (of Aff.)}\\
% City, Country \\
% email address or ORCID}
% \and
% \IEEEauthorblockN{3\textsuperscript{rd} Given Name Surname}
% \IEEEauthorblockA{\textit{dept. name of organization (of Aff.)} \\
% \textit{name of organization (of Aff.)}\\
% City, Country \\
% email address or ORCID}
% \and
% \IEEEauthorblockN{4\textsuperscript{th} Given Name Surname}
% \IEEEauthorblockA{\textit{dept. name of organization (of Aff.)} \\
% \textit{name of organization (of Aff.)}\\
% City, Country \\
% email address or ORCID}
% \and
% \IEEEauthorblockN{5\textsuperscript{th} Given Name Surname}
% \IEEEauthorblockA{\textit{dept. name of organization (of Aff.)} \\
% \textit{name of organization (of Aff.)}\\
% City, Country \\
% email address or ORCID}
% \and
% \IEEEauthorblockN{6\textsuperscript{th} Given Name Surname}
% \IEEEauthorblockA{\textit{dept. name of organization (of Aff.)} \\
% \textit{name of organization (of Aff.)}\\
% City, Country \\
% email address or ORCID}
% }

\maketitle

\begin{abstract}
Accurate forecasting of migration flows is critical for humanitarian aid but is challenged by sparse, short-history data. Current methods for selecting exogenous variables (features) rely on subjective domain knowledge or on data-driven approaches that fail in small data regimes. To address this, we propose a novel method that leverages the internal knowledge of Large Language Models (LLMs) to guide feature selection for time series forecasting. Our approach prompts an LLM to identify the most contextually relevant features from a candidate list, a process we enhance using Retrieval-Augmented Generation (RAG) to incorporate up-to-date, domain-specific knowledge. We evaluate our method on European migration data (FRONTEX, 2019–2024) and standard forecasting benchmarks. Our experiments show that LLM-guided feature selection consistently improves forecast accuracy over models using the full feature set and outperforms traditional data-driven techniques in data-scarce scenarios. Our work demonstrates the potential of LLMs for knowledge-driven feature selection in domains limited by data availability.
\end{abstract}

\begin{IEEEkeywords}
time series forecast, LLMs, feature selection, small data regime
\end{IEEEkeywords}

\section{Introduction}
Forecasting migration flows is a critical yet formidable challenge, essential for effective humanitarian resource allocation in response to escalating global displacement \cite{world_migration_report}. During crises such as the 2015–2016 European migration crisis, the failure of early warning systems highlighted the urgent need for reliable forecasts to guide timely aid \cite{mcmahon2021, napierala2022}. The complexity of this task is rooted in sparse, short-history data and the myriad of socio-economic factors driving human movement.

Incorporating exogenous factors as features into time series models can help by providing indirect information that influences human movement. These are often collectively referred to as push-pull factors---\textbf{push factors} motivate individuals to leave their home country (e.g., poverty, impact of bad weather season on agriculture), while \textbf{pull factors} attract them to a destination (e.g., job prospects, government stability)~\cite{lee1966theory}. %Ultimately, the decision to migrate is complex, and factors are weighted at an individual level.
Recently, machine learning models that use a relatively small number of push--pull factors were proposed for this task, achieving slight improvements over classic statistical models ~\cite{abel2013forecasting,  azizi2020artificial, fantazzini2021forecasting,  Carammia2022, qi2023modelling}. However, these results are not generalizable: the set of relevant factors differs across regions. Currently, the selection of exogenous variables is highly dependent on domain knowledge and therefore prone to human biases.

\begin{table}[t!]
\caption{Mean Absolute Percentage Error (MAPE) for `None' ($n=0$) and `All' ($n=157$) feature scenarios. Using all features causes overfitting; MAPE grows several folds.}
\label{tab:baseline_mape}
\begin{tabular}{lrrrrrr}
\toprule
Model & \multicolumn{2}{c}{Prophet} & \multicolumn{2}{c}{SARIMAX} & \multicolumn{2}{c}{TimeXer} \\
 & None & All & None & All & None & All \\
Route &  &  &  &  &  &  \\
\midrule
C. Medit. & \textbf{27.2} & 98.0 & \textbf{54.6} & 73.6 & \textbf{44.8} & 53.2 \\
E. Borders & \textbf{67.7} & 1050.5 & \textbf{726.4} & - & \textbf{71.6} & 74.9 \\
E. Medit. & \textbf{89.8} & 387.1 & \textbf{75.1} & no conv. & \textbf{117.8} & 315.5 \\
W. Medit. & \textbf{36.2} & \textbf{36.2} & 38.3 & \textbf{36.2} & 36.9 & \textbf{35.9} \\
\bottomrule
\end{tabular}
\end{table}

While it would be desirable to include a large number of potential factors to improve prediction accuracy, the fact that statistics on many of these variables are sparsely collected (e.g., monthly or quarterly) combined with the short-history of publicly available data related to migration may result in lack of data to train effective models from scratch.
%
Time Series Foundation Models (TSFMs) such as TimeXer emerge as a promising approach to leverage neural networks pre-trained on large and diverse datasets for few-shot learning~\cite{Liang_Wen_Nie_Jiang_Jin_Song_Pan_Wen_2024}. However, simple experiments shown in Table~\ref{tab:baseline_mape} and later elaborated in Section~\ref{sec:experiments} indicate that the error for both classic and deep time series models grows significantly when provided with too many features.

To select relevant task-dependent features for time series forecast while avoiding human biases, we turn our attention to Large Language Models (LLMs) and Retrieval Augmented Generation (RAG), which have consistently demonstrated the capability to capture contextual relationships and encode general knowledge across a variety of domains. This capability has been recently harnessed for feature selection in regression tasks~\cite{hegselmann2023tabllm, li_exploring_2025}. Inspired by these successes, we investigate whether these ideas can be adapted for predicting migration flows and, more generally, for time series forecast in small data regimes. Specifically, we set out to answer the following research questions:
\begin{enumerate}[label=\textbf{RQ\arabic*}]
    % \item \textbf{(Learning without feature selection):} are classic or deep time series models able to learn the relationship between large numbers of features and the target variable?
    \item \textbf{(Classic feature selection methods):} How well do current techniques for feature selection in forecast tasks work in small data regimes?
    \item \textbf{(LLM-based feature selection):} Can LLMs be used for feature selection and how they compare with existing techniques?
    \item \textbf{(Impact of LLM choice):} Does the choice of LLM impact the ability of selecting good features?
    \item \textbf{(Impact of RAG):} Can RAG be useful in guiding LLMs to select good features?
\end{enumerate}

\noindent \textbf{Our approach.} To investigate the previous research questions, we propose \method{}, a new \underline{LLM}-based method for \underline{T}ime \underline{S}eries \underline{F}eature \underline{S}election, especially designed for small data regimes. \method{} prompts an LLM to select contextually-appropriate features from a candidate list for a given task, which are then used for training and/or inference with time series forecasting models. \fm{Can we add more details on how it works?} We conduct a comprehensive array of experiments on benchmark datasets to characterize the effectiveness of different LLMs when choosing $k$ features for different values of $k$, considering all combinations of two LLMs (Gemini-2.5 and GPT4o-mini) and four time series models (SARIMAX, Prophet, TimeXer and iTransformer).

%
Moreover, we investigate the impact of adding a Retrieval Augmented Generation (RAG) module, which enables the use of documents to incorporate up-to-date domain-specific knowledge into the feature selection. For this investigation, we conduct a case study on migration flows into the European Union between 2019 and 2024 using a dataset collected from FRONTEX immigration data. For the exogenous variables, we collect XXX factors from source countries including GDP, \fm{include others}. In addition, we collect a corpus of XXX documents from Wikipedia? related to the human flows through the five main routes into the EU.


In summary, here we make four main contributions:
\begin{itemize}
    \item We explore, for the fist time in the literature, the feasibility of using TSFMs for forecasting migration flows based on a large number of push-pull factors.
    \item We propose the \method{} method, which pioneers the use of LLMs for feature selection in time series forecast tasks.
    \item We characterize the performance of \method{} on various benchmark datasets under various data regimes.
    \item We conduct a case study on the FRONTEX immigration data and evaluate the impact of RAG on \method{}'s performance. 
\end{itemize}

\fm{Add findings here.}

\section{Related Works}

\subsection{Adapting LLMs for Time Series Prediction} Large Language Models (LLMs) have demonstrated a remarkable ability to model dependencies across sequential data, which has led to interest in adapting them to time series forecasting task. However, directly applying LLMs to this task is challenging because time series are typically described by numerical values with strong temporal dependencies. Instead of training from scratch, most approaches preserve LLM architecture and weights, modifying the input format of time series data to match the LLM's expectation. Prominent approaches range from directly encoding time series data as strings of digits and framing the task as next token prediction~\cite{gruver2024largelanguagemodelszeroshot} to segmenting time series into patches that are passed onto transformer encoders~\cite{li2024llm4ts,zhou2023fitsallpowergeneraltime, jin2024timellmtimeseriesforecasting}. LLM4TS \cite{li2024llm4ts} incorporates instance normalization and channel independence before patching, while Time-LLM \cite{jin2024timellmtimeseriesforecasting} introduces ``patch reprogramming'', utilizing both word and patch embeddings through a multi-head cross-attention mechanism.
Beyond input adaption, researchers have also explored fine-tuning strategies to enable LLMs to perform forecasting tasks effectively without losing their pre-trained knowledge. GPT4TS \cite{zhou2023fitsallpowergeneraltime} experimented with freezing the attention mechanism while fine-tuning only the input and positional embeddings. LLM4TS \cite{li2024llm4ts} offers a more intricate/introduces a two-stage fine-tuning approach that includes a specialized time series alignment phase with multiple encoding layers. This evolving landscape of techniques demonstrates an effort to bridge the gap between the architectural strengths of LLMs and the dependent characteristics of time series data./Nevertheless, none of the current techniques leverages the LLM's internal knowledge about the semantics of the time series variables for prediction.

\subsection{Time Series Foundational Models (TSFM)}
Time Series Foundation Models (TSFM) represent a paradigm shift in forecasting, enabling pre-trained models to generalize across diverse temporal datasets \cite{Liang_Wen_Nie_Jiang_Jin_Song_Pan_Wen_2024}. However, incorporating exogenous variables remains challenging, with most foundation models designed primarily for univariate forecasting \cite{Liang_Wen_Nie_Jiang_Jin_Song_Pan_Wen_2024}. TimeXer~\cite{TimeXerNEURIPS2024} represents the first transformer architecture explicitly designed for exogenous variable integration, introducing a dual-attention mechanism that employs patch-wise self-attention for endogenous variables and variate-wise cross-attention for exogenous interactions. Salesforce's Moirai \cite{WOO_LIU_KUMAR_XIONG_SAVARESE_SAHOO_2024} addresses this through Any-Variate Attention, handling arbitrary variates via flattened sequences with binary attention biases, trained on 27 billion observations. IBM's Tiny Time Mixers (TTM) \cite{Ekambaram_Jati_Dayama_Mukherjee_Nguyen_Gifford_Reddy_Kalagnanam_2024} challenge transformer dominance using MLP-based mixing with dedicated exogenous blocks, achieving competitive performance with $<1$M parameters. Current fusion strategies include early fusion (concatenation), late fusion (separate processing with cross-attention), and hybrid approaches \cite{TimeXerNEURIPS2024, Ekambaram_Jati_Dayama_Mukherjee_Nguyen_Gifford_Reddy_Kalagnanam_2024}. Recent work on CITRAS \cite{Yamaguchi_Suemitsu_Wei_2025} and ExoTST \cite{Tayal_Renganathan_Jia_Kumar_Lu_2024} demonstrates up to 10\% accuracy improvements through proper temporal alignment of past and current exogenous variables. Critical limitations persist, including limited native exogenous support in most models, temporal misalignment challenges, and scalability-interpretability trade-offs \cite{Liang_Wen_Nie_Jiang_Jin_Song_Pan_Wen_2024}. 

\subsubsection{Leveraging LLMs for Feature Selection.} 
There are several works exploring LLMs for feature selection. 
LMPriors pioneered using LLMs for feature selection by instructing GPT-3 to mark importance of features ("y"/"n") based on feature names and descriptions, selecting features if the difference in log-probabilities for generating "y" and "n" tokens reached a threshold\cite{choi_lmpriors_2022}. LLM-select builds on this by proposing three pipelines that rely solely on textual information and directly using the output of generated text \cite{jeong_llm-select_2024}. ICE-SEARCH further refines feature selection by prompting a LLM to filter features based on an externally computed score \cite{yang_ice-search_2024}, while FeatLLM advances feature engineering by incorporating conditional rules to generate additional meta-features to improve predictions \cite{han_large_2024}. FREEFORM integrates chain-of-thought and emsembling features to select and engineer based on the knowledge of LLMs for several genotype datasets \cite{lee_knowledge-driven_2024}. Recently, research has shifted into integrating LLMs for feature selecti
on with more traditional data-driven methods. For instance, LLM-Lasso~\cite{zhang_llm-lasso_2025} uses LLMs to generate penalty factors for each feature that are then converted into weight factors for the Lasso penalty, while LLM4FS~\cite{li_llm4fs_2025} integrates contextual knowledge from LLMs with traditional data-driven techniques such as random forest and forward sequential selection. These works illustrate the potential of LLMs to encode relevant, contextual information about a given task to augment traditional, data-driven, supervised learning approaches. However, the present study is the first to explore this idea in the context of time series forecast.
%\fm{Currently this subsection of RW is a bit long. You can try to combine some works in a single sentence, e.g., LLM-Lasso and LLM4FS.}
%There are several works exploring LLMs for feature selection. 
%\citet{choi_lmpriors_2022} pioneered this work in their framework LMPriors, where they instruct GPT-3 to mark each feature as important (“Y”) or unimportant (“N”) based on feature names and descriptions and select features if the difference in log-probabilities for generating “Y” and “N” tokens reaches a threshold. \fm{Does the next technique have a name? I'm leaning towards referring to the works by the names and then using the cite command.} Building upon this, \citet{jeong_llm-select_2024} proposed three pipelines that rely solely on textual information and directly use the output of generated text. Additional studies explore more complex pipelines for feature selection and feature engineering \cite{yang_ice-search_2024, han_large_2024, lee_knowledge-driven_2024}. \citet{yang_ice-search_2024} introduce In Context Evolutionary Search, ICE-SEARCH, which iteratively optimizes feature selection by prompting an LLM to filter features based on externally computed test scores. \citet{han_large_2024} use LLMs as feature engineers through their FeatLLM framework, where conditional rules generate additional meta-features to improve predictions. \citet{lee_knowledge-driven_2024} introduce FREEFORM which integrates chain-of-thought and ensembling features to select and engineer features based on the knowledge of LLMs for several genotype datasets. More recently, work has been done to integrate LLMs for feature selection with more traditional data-driven methods. For instance, LLM-Lasso~\cite{zhang_llm-lasso_2025} uses LLMs to generate penalty factors for each feature that are then converted into weight factors for the Lasso penalty, while LLM4FS~\cite{li_llm4fs_2025} integrates contextual knowledge from LLMs with traditional data-driven techniques such as random forest and forward sequential selection. These works illustrate the potential of LLMs to encode relevant, contextual information about a given task to augment traditional, data-driven, supervised learning approaches. \fm{However, the present study is the first to explore this idea in the context of time series forecast.}
% \\\\
% \subsubsection{Machine Learning Approaches for Forced Migration Flows}
% \begin{itemize}
%     \item broadly introduce the growing interest/application of ML to forced migration
%     \item main types of ML models/techniques - predictive?/descriptive?/analytical? - probs predictive only
%     \item specific problems of forced migration prediction that ML approaches try to solve
%     \item strengths/limitations of these approaches

% \end{itemize}

\section{Proposed Method}

\begin{figure}
    \centering
    \includegraphics[width=0.9\linewidth]{figs/overview_fig_llm.pdf}
    \caption{PLACEHOLDER; figure needs to be updated}
    \label{fig:placeholder}
\end{figure}

We employ Large Language Models (LLMs) as structured knowledge repositories to guide intelligent feature selection, enhanced by Retrieval-Augmented Generation (RAG) techniques, for domain specific contexts. This approach operates in four stages:
% (RAG Database preparation): collect a corpus of documents that may provide information about the relationship between exogeneous features and the target time series.

% (Optional RAG step): build a query from a template use it to query the RAG database, then use the retrieved documents as part of the context for the next step

% (LLM-based ranking of candidate features): \fm{DESCRIBE THE LOGIC USED FOR BUILDING THE TEMPLATE. WHAT ARE THE MAIN SECTIONS? EMPHASIS IN THE SELECTION GUIDELINES} we ask it to select the top $k=10$ (or 20?) features and rank it; we repeat the process N=3 times to account for randomness and use the outputs as an ensemble model; we take the output and run a deterministic algorithm to extract the text and select the features. then we sort by the number of occurrences -- we save these intermediate results (this allows us later to pass the same features to the different time series models) -- and cut it based on the number of features that we actually use; \fm{if we have time, add an algorithm env}



% (Training of time series model):
\begin{enumerate}
    \item \textbf{Semantic Analysis:} Given a set of candidate exogenous variables and a target variable, the LLM analyzes potential causal relationships and domain-specific relevance. \fm{merge:} \textbf{Structured Selection:} The LLM outputs a ranked selection of features with explicitly justification based on the causal relevance, and predictive power.
    \item \textbf{RAG Enhancement:} For domain specific datasets (e.g. migration), we augment the LLM's reasoning with retrieved contextual documents from a curated knowledge base, providing additional domain expertise.
\end{enumerate}

\textbf{RAG Database Preparation}: We begin by curating a knowledge base of relevant migration-related documents, such as reports from Frontex and other humanitarian organizations. These documents are segmented into overlapping text chunks using a recursive splitting strategy designed to preserve semantic coherence. Each chunk is embedded using the multilingual-E5-large-instruct model~\cite{wang2024multilingual}, which supports cross-lingual semantic similarity. The resulting embeddings are stored in a persistent vector database using ChromaDB, enabling efficient similarity search. 

ADD THIS: cleaning watermarks; recursive indexing; two-stage retrieval, BAAI models, the second uses cross-encoder for re-ranking;

\noindent\textbf{(Optional) RAG-Enhanced LLM Selection}: To support domain-specific reasoning, particularly for migration datasets, we enhance the LLM-based feature selection process with Retrieval-Augmented Generation (RAG). At inference time, a query derived from the forecasting task or target variable is used to retrieve the top-k most relevant document chunks. After that, the top-k results are reranked using a cross-encoder model. The retrieved content is then incorporated into the LLM prompt as contextual grounding, enabling the model to make more informed and context-aware feature selections based on external migration knowledge.\\


\section{Experimental Settings}\label{sec:experiments}

%\fm{This section describes the two approaches we introduce for leveraging LLMs' internal knowledge in time series forecast tasks and how RAG can be used to enhance them in specific domains.}


% \subsection{Approach 2: LLM-embedding Adapter for Time Series Forecast}

% We introduce context-aware embedding mechanisms that explicitly encode semantic relationships within Time Series Foundation Models (TSFMs), enabling dynamic behavioral modulation based on contextual understanding of variable interdependencies.

We perform a comprehensive set of experiments to address the Research Questions posed in Section~\label{sec:intro}. This section describes the datasets, time series models, feature selection baselines and the evaluation protocol.

\subsection{Data}

\textbf{FRONTEX migration data}: DESCRIBE DATA HERE
% \textbf{Migration}: sources, and what it is - add mods and credit to ARCHES in the ACK of the camera-ready version\\

Additionally, we use three popular time series benchmarks to assess the generalizability of our approach. Since we focus on small data regimes, we consider a subinterval of the original dataset's time span to restrict the available for for training the models, as described below.

% fix with elizabeth
\begin{itemize}
    \item \textbf{Weather}: sources, and what it is - add problematic same series \cite{weather_dataset}
    \item \textbf{Electricity}: sources, and what it is \cite{electricity_dataset}
    \item \textbf{ETTh/ETTm}: sources, and what it is - add how we modified \cite{ett_dataset}\\
\end{itemize}
\fm{PHONG: Add info about the collection of docs for RAG database}

\textbf{Retrieval corpus for RAG.}
We curate an external corpus to ground LLM reasoning: FRONTEX risk and route reports, WOLA policy briefs/analyses, Wikipedia country profiles (programmatically collected via ISO codes). Documents are cleaned, chunked, embedded to a vector database for using.

\subsection{Large Language Models}

Our core approach uses Large Language Models to identify semantically-relevant exogenous features based on domain knowledge and causal reasoning. We evaluate two efficient and commercial foundation models: GPT-4o-mini and Gemini Flash 2.0, chosen for their multilingual capabilities and fast inference speeds. Feature selection is then chosen by a structured prompt template.
The prompt include:
\begin{itemize}
\item Role: Define the LLM's responsibility in the forecasting process.
\item Background: Describe the forecasting goal and dataset context
\item Task: Select the top-k most impactful exogenous features.
\item Criteria: Causal relevance, predictive power, and lead/lag relationship.
\end{itemize}

\subsection{Models}

We selected a diverse set of forecasting models to comprehensively evaluate our LLM-based feature selection across different modeling paradigms and architectural approaches.

\begin{itemize}
\item \textbf{Seasonal ARIMA w/ exogenous factors (SARIMAX):} We include this as our primary statistical baseline due to its explicit support for exogenous variables and proven effectiveness in capturing seasonal patterns.

\item \textbf{Prophet:} Selected for its robust handling of trend changes and holiday effects, Prophet complements SARIMAX by offering a decomposable approach that can reveal how LLM-selected features interact with different time series components (trend, seasonality, holidays). Its flexibility with irregular patterns makes it particularly valuable for our diverse dataset collection.

\item \textbf{TimeXer:} A transformer-based model with efficient patching mechanisms for processing long sequences.

\item \textbf{iTransformer:} Features an inverted attention approach that treats variables as tokens rather than time steps.

\end{itemize}


% \textbf{TimesNet:} A CNN-based architecture that incorporates FFT-guided period detection for frequency-domain analysis.

% \textbf{PatchTST:} A patch-based transformer that serves as both a competitive neural baseline and, in univariate mode, a comparison point against our multivariate feature-enhanced approaches.

\subsection{Feature Selector Baselines}

To comprehensively evaluate our approach, we compare LLM-based feature selection against established baselines and ablation conditions:
\begin{itemize}
\item \textbf{No Features}: Pure autoregressive forecasting using only historical target values, establishing the baseline performance without external information.
\item \textbf{All Features}: Includes every available exogenous variable, representing the upper bound of available information while potentially introducing noise and overfitting.
\item \textbf{XGBoost Selection}: Gradient boosting-based feature importance ranking, representing a current state-of-the-art method in automated feature selection for structured data.
% \item \textbf{Random Selection}: Randomly sampled features matching the number selected by LLM methods, controlling for feature set size effects and establishing a lower-bound baseline.
\end{itemize}






with dataset-specific forecast horizons optimized for practical applications.


[add text for how we perform our experiments - expanding window/ trainging etc etc]

\subsection{Training and Inference}

\fm{PHONG: could you add info about the sliding window and size of initial windows, forecast horizon, step size}

For datasets where the target time series is non-negative, we clip the predictions: $\widehat{y}_t = \max(F(\mathbf{x}_{t-\tau:t}), 0)$.


\subsection{Evaluation}



We evaluate the prediction using standard regression metrics: Root Mean Square Error (RMSE), Mean Absolute Percentage Error (MAPE). In each case, we average the predictions across all forecasts (see sliding window in Section~\ref{sec:}). We consider a smoothed version of MAPE to prevent division by small numbers:
\begin{equation}\label{eq:smooth_mape}
    \mathrm{MAPE}(\mathbf{\widehat{y}}_{t:t+h},\mathbf{y}_{t:t+h}) = \frac{1}{h}\sum_{i=0}^{h-1}\frac{|\widehat{y}_{t+i} - y_{t+i}|}{|y_{t+i}+10|}.
\end{equation}

When conducting a more holistic evaluation over the entire dataset, we do not take averages of error metrics across different routes because the ground-truth values vary significantly in magnitude. Instead, we rank the performance of different configurations for each route.

Formally, we define a configuration $c = (m, f, n)$ as a tuple where $m$ is a model, $f$ is a feature selector and $n$ is the number of features. Each configuration $c$ is associated a set of performance metrics $h(c) \in \mathbb{R}^d$. We can rank a set of configurations $\mathcal{C}$ by sorting them according to a chosen performance metric (i.e., a dimension in $h(c)$). The best performing configuration has rank 1, while the worst has rank $|\mathcal{C}|$.

We aggregate the performance across routes by computing the mean reciprocal rank (MRR):
\begin{equation}\label{eq:mrr}
    \mathrm{MRR}_\mathrm{metric}(c,\mathcal{C}) = \frac{1}{N}\sum_{r=1}^{R} \frac{1}{\mathrm{rank}_\mathrm{metric}(c,\mathcal{C},r)},
\end{equation}
where $R$ is the number of routes and $\mathrm{rank}_\mathrm{metric}(c,\mathcal{C},r)$ is rank of configuration $c$ among the set of  on the route indexed by $r$.

\section{Results}

We conduct preliminary experiments to investigate whether time series models can be trained effectively when a large number of features (\texttt{All}) are provided as input to the model. In Table~\ref{tab:baseline_mape} (Section~\ref{sec:intro}), we contrast each model's performance against the setting where no features are provided (\texttt{None}). While using all features led to a sizable reduction in error for Prophet on Central Mediterranean, in other cases the benefit was relatively small or the error has gone up by several folds. TimeXer showed the smallest degradation in comparison to the reference setting, except for the Western Balkan route.

To understand whether including some intelligently-selected features can outperform the best among the two baselines, we conduct a large set of experiments varying the number of features in  $\mathcal{N} = \{1,2,3,5,10,15\}$ for each pair of model and feature selector. We focus on the three routes where the relative performance difference between \texttt{None} and \texttt{All} is the largest: Central, Eastern and Western Mediterranean.

\begin{figure}
    \centering
    \includegraphics[width=0.99\linewidth]{figs/plot_RMSE_by_Model.pdf}
    \caption{Caption}
    \label{fig:placeholder}
\end{figure}

Next, we evaluate \method{} (\textbf{RQ2}) and compare it against a traditional feature selection method (\textbf{RQ1}). We choose the XGBoost Selection for being a popular choice for feature selection in the context of time series forecast. Rather than just finding out which configuration performs best, our goal is to understand the impact of feature selection on each model individually. Therefore, we compute each model's overall performance in terms of the Mean Reciprocal Rank (MRR), defined in Eq.~\ref{eq:mrr}, by ranking each configuration $c=(m,f,n)$ relative to the set $\mathcal{C} = \{(m, f',n'): f' \in \mathcal{F}, n'\in \mathcal{N'} \}$ of all configurations that share the same time series model $m$, where $\mathcal{F}$ is the set of feature selectors and $\mathcal{N} = \{1,2,3,5,10,15\}$ contains options for number of features.



\begin{figure}
    \centering
    \includegraphics[width=0.99\columnwidth]{figs/plot_mrr_ranking_selector_features_by_rmse.pdf}
    \caption{Replace by PDF figure}
    \label{fig:mrr}
\end{figure}

Figure~\ref{fig:mrr} shows the MRR for each configuration $c = (m,f,n)$ grouped by model $m$. For Prophet, no configuration has a particularly high MRR, but \method{} tends to perform well when selecting a small number of features ($n=1$, 2 or 3). XGBoost performs best for $n=5$ and $15$. SARIMAX, in turn, benefits from having few features ($n=1$ or 2). In that setting, feature selected by \method{} are much more effective than the ones selected by XGBoost. WHY IS THAT? IT LOOKS LIKE THE FEATURES INITIALLY SELECTED BY XGBOOST ARE NOT EASILY HARNESSED BY SARIMAX.

\fm{Maybe leave TimeXer out of this figure. I think the other two have a more interesting behavior}. Based on Figure~\ref{fig:mrr}, we observe that using Gemini as feature selector matches or exceeds the performance or GPT4o-mini, except when TimeXer is used for forecast.



\section{Case Study using RAG}

\section{Discussion}

\section{Conclusion}

We introduce a novel method to combine machine learning approaches in time series forecasting with the power of LLMs. Our work leverages LLMs to perform feature selection of exogenous variables for domain-specific contexts. Our method prompts an LLM to select the most influential and relevant features for time series forecasting in a particular domain from a list of possible variables, and then performs time series forecasting with the selected variables. This process can be enhanced by Retrieval Augmented Generation. 

\section{Limitations \& Future Work}


\bigskip
\setcounter{secnumdepth}{1}
\appendix

\section{Prompt Templates}\label{app:prompts}
Here we provide the full prompt for each dataset.
\section{RAG Implementation Details}\label{app:rag}
\textbf{Corpus:} Frontex reports, ISOm etc. \\
\textbf{Chunking:} recursive, size=800, overlap=200. \\
\textbf{Embeddings:} multilingual-e5-large-instruct (dim=1024). \\
\textbf{Index:} ChromaDB; cosine similarity; top-$k=5$. \\
\textbf{Prompt integration:} .


\definecolor{appendixpurple}{RGB}{152, 78, 163}

\begin{tcolorbox}[breakable, colback=appendixpurple!10!white, colframe=appendixpurple!50!black, title=LLM-as-Judge Prompt]
\textbf{System Prompt:} You are an expert assistant for time series forecasting.
Your primary objective is to select the most relevant exogenous features to improve forecasting accuracy.

\textbf{Background:}  
We are working with the \{dataset\_name\} dataset, \{dataset\_description\}.  
The target is \{target\}.

\textbf{Task:}  
From the list below, choose the top \{n\_features\} most impactful exogenous variables to include as inputs when forecasting \{target\}.

\textbf{Selection Guidelines:}  
1) Causal influence (highest priority): Does the exogenous variable meaningfully affect the \{target\}?  
2) Predictive power: Does it show strong correlation or predictive relationship with the \{target\}?

Prioritize retrieved context documents to support your decisions. If none are available, rely on domain knowledge.

\textbf{Formatting Rules:}  
You must \textbf{ONLY} return your final answer between the following special markers — nothing else should be printed.

\begin{verbatim}
BEGIN FINAL ANSWER SECTION
@@START_RANKED_FEATURES@@
<feature_name>: concise justification
<feature_name>: concise justification
...
@@END_RANKED_FEATURES@@
END FINAL ANSWER SECTION
\end{verbatim}

The list of candidate variables is:

\end{tcolorbox}





% this ended up working? If we must convert, let me know and I will do it. SMS
\bibliographystyle{IEEEtran}
\bibliography{references}

\pagebreak
\onecolumn
\include{full_table}
\include{full_table_new}


\end{document}
